%% kapa: exam=1 round=20 year=2009 period=1 form=A subject=economics subject_ko=경제원론 seq=15 number=55
%% source: Q-Net 감정평가사 기출문제 (cst00309) artlSeq=5024178 file=extracted/1008/5024178/2009년 제20회 감정평가사 1차시험 기출문제/1교시.hwp
%% conversion: kordoc 4.24.51 -> tools/gen1.py (auto); review=proofread-sonnet; flags=none
\begin{question}{55}
두 개의 분리된 시장에서 동일한 상품을 판매하는 독점기업이 있다. 두 시장의 수요곡선은 각각 $P_{ 1 } =100-0.5Q_{ 1 }$ 과 $P_{ 2 } =80-Q_{ 2 }$ 이다. 각 시장에 대해 가격차별이 가능하고, 소비자 간 전매는 불가능하다. 이윤을 극대화하는 독점기업이 더 높은 가격을 설정하는 시장과 두 시장 간의 가격 차는? (단, 공급자의 한계비용은 0이다. $P_{ 1 }$ 과 $P_{ 2 }$ 는 각각 시장 1과 시장 2의 상품가격이며, $Q_{ 1 }$ 과 $Q_{ 2 }$ 는 각각 시장 1과 시장 2의 수요량이다.)
\begin{choices}
  \item 시장 1, 가격 차 5
  \item 시장 1, 가격 차 10
  \item 시장 2, 가격 차 5
  \item 시장 2, 가격 차 10
  \item 가격차별을 하지 않는다.
\end{choices}
\end{question}
\answer{2}
